% Copyright (c) 2025 Intel Corporation.  All rights reserved.  See the file COPYRIGHT for more information.
% SPDX-License-Identifier: Apache-2.0

\documentclass[11pt]{article}

% Page layout
\usepackage[letterpaper, margin=1in]{geometry}
\setlength{\headheight}{14pt}
\usepackage{parskip}
\setlength{\parindent}{0pt}

% Typography
\usepackage{fancyvrb}
\usepackage{booktabs}
\usepackage{amsmath}
\usepackage{amssymb}

% Code listings
\usepackage{listings}
\lstset{
    basicstyle=\ttfamily\small,
    columns=fullflexible,
    keepspaces=true,
    showstringspaces=false,
    breaklines=true,
    frame=none,
    xleftmargin=2em,
}

% Hyperlinks
\usepackage[colorlinks=true, linkcolor=blue, urlcolor=blue, citecolor=blue]{hyperref}

% Headers and footers
\usepackage{fancyhdr}
\pagestyle{fancy}
\fancyhf{}
\fancyhead[L]{\textit{Formal Verification for the Fulcrum CSP Toolchain}}
\fancyhead[R]{\thepage}
\renewcommand{\headrulewidth}{0.4pt}

\title{\textbf{Formal Verification for the Fulcrum CSP Toolchain} \\[0.5em]
\large Feasibility Study, Project Plan, and Implementation Report}
\author{Mika Nystr\"om \and Claude (Anthropic)}
\date{DRAFT Version 0.2\\
\today}

\begin{document}

\maketitle

\vspace{1em}

\noindent
 Copyright \copyright\ 2025 Intel Corporation

 Licensed under the Apache License, Version 2.0 (the ``License'');
 you may not use this file except in compliance with the License.
 You may obtain a copy of the License at

 {\tt http://www.apache.org/licenses/LICENSE-2.0}

\tableofcontents
\newpage

%======================================================================
\section{Introduction}
%======================================================================

Asynchronous circuit design offers compelling advantages in power
efficiency, modularity, and electromagnetic compatibility, but it
introduces verification challenges that synchronous flows do not face.
The absence of a global clock means that correctness depends entirely on
the interplay of handshaking protocols: a single missing acknowledgment
or misordered communication can cause a physical deadlock that is
invisible to simulation until the right interleaving is exercised.

The Fulcrum CSP toolchain compiles Communicating Sequential Processes
(CSP) descriptions into cycle-accurate Modula-3 simulators.  Simulation
is effective for validating functional behavior, but it explores only a
subset of the reachable state space.  Formal verification---exhaustive
or symbolic exploration of \emph{all} reachable states---can provide
mathematical guarantees of deadlock freedom, livelock freedom, and
protocol conformance that simulation alone cannot.

This report assesses the feasibility of adding FDR4-style formal
verification capabilities to the Fulcrum CSP toolchain.  It surveys the
relevant theory, compares the Fulcrum CSP dialect with the CSPM input
language of FDR4, identifies favorable structural properties and
key challenges, proposes three implementation approaches, and
recommends an incremental path forward with a phased project plan.

%======================================================================
\section{Background}
%======================================================================

%----------------------------------------------------------------------
\subsection{CSP Theory and Semantic Models}
%----------------------------------------------------------------------

CSP (Communicating Sequential Processes), introduced by Hoare~\cite{hoare1985},
provides a process algebra for describing systems of concurrent agents
that interact by synchronous message passing over named channels.
The theory defines three progressively richer semantic models:

\begin{description}
\item[Traces model ($\mathcal{T}$).]
  A process is characterized by the set of finite sequences of
  events (traces) it can perform.  The traces model supports
  safety properties: ``nothing bad happens.''

\item[Stable failures model ($\mathcal{F}$).]
  A process is characterized by its traces together with its
  \emph{failures}---pairs $(s, X)$ where $s$ is a trace and $X$ is a
  set of events the process can refuse after performing~$s$.
  The failures model supports deadlock analysis and determinism
  checking.

\item[Failures-divergences model ($\mathcal{FD}$).]
  Extends the failures model with \emph{divergences}---traces after
  which the process can perform an infinite sequence of internal
  ($\tau$) actions.  This is the richest standard model and supports
  livelock analysis.
\end{description}

Refinement checking is the central verification technique: a system
$\mathit{Impl}$ satisfies a specification $\mathit{Spec}$ if
$\mathit{Impl}$ refines $\mathit{Spec}$ in the chosen model, written
$\mathit{Spec} \sqsubseteq \mathit{Impl}$.  In the traces model this
means the implementation produces no trace that the specification
forbids; in the failures model it additionally means the implementation
refuses no more than the specification allows.

%----------------------------------------------------------------------
\subsection{FDR4: Architecture and Algorithms}
%----------------------------------------------------------------------

FDR4~\cite{fdr4} (Failures-Divergences Refinement checker, version~4)
is the standard tool for CSP refinement checking.  Its architecture
comprises several stages:

\begin{enumerate}
\item \textbf{Parsing.}  The input language is CSPM (machine-readable
  CSP), a functional language with pattern matching, let-bindings,
  comprehensions, and built-in types (sets, sequences, tuples).
  Processes are defined using CSP operators: prefix ($\to$),
  external choice ($\Box$), internal choice ($\sqcap$), parallel
  composition ($\|$), interleaving ($|||$), hiding, and renaming.

\item \textbf{Supercombinator compilation.}  Process expressions are
  compiled into \emph{supercombinators}---compact representations that
  enumerate the transitions of a process as a function of its
  parameters.  This is the key innovation of FDR4: supercombinators
  enable on-the-fly state enumeration without materializing the full
  state graph up front.

\item \textbf{State exploration.}  FDR4 explores the product state
  space using BFS or DFS, driven by the supercombinator transition
  functions.  Compression techniques (diamond, sbisimulation,
  $\tau$-loop elimination) reduce the state space during exploration.

\item \textbf{Refinement checking.}  The tool checks
  $\mathit{Spec} \sqsubseteq_X \mathit{Impl}$ where $X$ is the
  chosen semantic model ($\mathcal{T}$, $\mathcal{F}$, or
  $\mathcal{FD}$).  Special-case checkers exist for deadlock freedom,
  livelock freedom, and determinism, which do not require an explicit
  specification process.
\end{enumerate}

FDR4 achieves scalability through several mechanisms: symbolic
representation of transitions, compression of the explored state
space, and multi-core parallel exploration.  Nevertheless, state
explosion remains the fundamental limit: a system of $n$ processes
each with $k$ states has a product space of up to $k^n$ states.

%======================================================================
\section{The Fulcrum CSP Toolchain}
%======================================================================

The Fulcrum CSP toolchain compiles hardware-oriented CSP descriptions
into native simulator binaries.  Its architecture is:

\begin{Verbatim}[samepage=true]
  .csp source          .sys system description
       |                        |
       v                        v
    cspfe (parser)       sys-file parser (Scheme)
       |                        |
       v                        v
  S-expression AST       system graph (channels,
  (.scm files)           process instances, ports)
       |                        |
       +-------+    +-----------+
               |    |
               v    v
     cspc (9-pass Scheme compiler)
               |
               v
     Modula-3 state machines
     (m3__*.m3, SimMain.m3, m3makefile)
               |
               v
     cm3 -build -override
               |
               v
     Native simulator binary (sim)
\end{Verbatim}

\subsection{Key Components}

\begin{description}
\item[\texttt{cspparse} (\texttt{cspfe}).]
  Lex/Yacc-based parser that reads \texttt{.csp} source files and
  emits an S-expression representation of the AST.  The AST is defined
  by the \texttt{cspgrammar} library.

\item[\texttt{cspgrammar}.]
  Modula-3 library defining the abstract syntax tree.  Key types
  include \texttt{CspExpression} (literals, binary/unary operators,
  channel operations like \texttt{Receive}, \texttt{Peek},
  \texttt{Probe}), \texttt{CspStatement} (assignment,
  \texttt{DetSelection}, \texttt{NondetSelection},
  \texttt{DetRepetition}, \texttt{Sequential}, \texttt{Parallel},
  \texttt{Send}, \texttt{Recv}), and \texttt{CspType} (integers,
  booleans, channels, arrays, structures).

\item[\texttt{cspc}.]
  The compiler driver.  Invoked with \texttt{-scm} it enters a Scheme
  REPL (via the embedded \texttt{mscheme} interpreter) that drives a
  9-pass compilation pipeline.  Scheme stubs expose the Modula-3 AST
  types to Scheme code, which performs type checking, optimization,
  and code generation.

\item[\texttt{cspsimlib}.]
  Modula-3 simulation runtime providing the scheduler, channel
  implementations (with configurable slack), process dispatch,
  and multi-threaded worker coordination.

\item[\texttt{cspbuild}.]
  A Scheme-level build system that reads \texttt{.sys} files
  (multi-process system descriptions), orchestrates compilation of
  each process, and generates the final simulator binary.
\end{description}

\subsection{Simulation Capabilities}

The generated simulator supports multiple scheduling strategies
(\texttt{-greedy}, \texttt{-nondet}, \texttt{-eager}), multi-threaded
execution (\texttt{-mt}~$N$), and distributed simulation
(\texttt{-master}/\texttt{-worker}).  It exercises one execution
path per run; randomized scheduling (\texttt{-nondet}) can improve
coverage but cannot guarantee exhaustive exploration.

%======================================================================
\section{Language Comparison: Fulcrum CSP vs.\ CSPM}
%======================================================================

The Fulcrum CSP dialect (sometimes called CAST) is designed for
hardware process descriptions, while CSPM is a general-purpose
process algebra notation.  The following table summarizes the key
differences:

\begin{center}
\begin{tabular}{lll}
\toprule
\textbf{Feature} & \textbf{Fulcrum CSP / CAST} & \textbf{CSPM (FDR4)} \\
\midrule
Channel model     & Bundled-data, directional & Abstract events \\
                  & (\texttt{L?x}, \texttt{R!x}) & (\texttt{c.v -> P}) \\
Data types        & Fixed-width integers, arrays, & Sets, sequences, tuples, \\
                  & structs & user-defined datatypes \\
Choice            & Deterministic (\texttt{[\,]}) & External ($\Box$), \\
                  & nondeterministic (\texttt{:[\,]:}) & internal ($\sqcap$) \\
Repetition        & \texttt{*[\,]} (infinite loop) & Recursion \\
Parallelism       & \texttt{||} (interleaved), & $\|$, $|||$, $[|\ldots|]$ \\
                  & implicit via \texttt{.sys} & with alphabet control \\
Hiding            & Not explicit & \texttt{P \textbackslash\ X} \\
Renaming          & Port bindings in \texttt{.sys} & \texttt{P[[a<-b]]} \\
Specification     & None (simulation only) & Refinement assertions \\
State space       & Implicit (simulated) & Explicit (enumerated) \\
\bottomrule
\end{tabular}
\end{center}

Despite the syntactic differences, the underlying computational model
is similar: both describe networks of sequential processes that
communicate via synchronous handshakes on named channels.  The Intel
dialect is closer to the ``blocking point'' formulation where a process
alternates between computation and a single blocking communication,
which is actually \emph{favorable} for verification---it constrains
the branching factor of the state graph.

%======================================================================
\section{Feasibility Analysis}
%======================================================================

%----------------------------------------------------------------------
\subsection{Favorable Factors}
%----------------------------------------------------------------------

Several properties of the Fulcrum CSP toolchain and the circuits it
describes make formal verification feasible:

\begin{description}
\item[Static topology.]
  The \texttt{.sys} file declares all processes and channels
  statically.  There is no dynamic process creation or channel
  allocation.  The system graph is known at compile time, which
  is a prerequisite for state enumeration.

\item[Bounded data.]
  Hardware processes operate on fixed-width integers (typically
  8--64~bits).  While the raw data domain is large, abstract
  interpretation or data-independence arguments can often reduce
  the effective data range.  Many protocol properties (deadlock
  freedom, channel protocol conformance) are independent of data
  values entirely.

\item[Blocking-point structure.]
  Fulcrum CSP processes alternate between local computation and a
  single blocking communication (send or receive on a channel).
  Between blocking points, the process executes deterministically.
  This means the LTS (Labelled Transition System) of each process
  has transitions only at communication points, dramatically reducing
  the number of states compared to a fine-grained interleaving
  semantics.

\item[Existing AST infrastructure.]
  The \texttt{cspgrammar} library provides a complete, well-typed
  AST.  The Scheme-accessible AST means that analysis passes can be
  prototyped rapidly in the existing \texttt{cspc} environment.

\item[Guarded-command normal form.]
  The \texttt{DetSelection} and \texttt{NondetSelection} statement
  types directly correspond to CSP external and internal choice.
  Guards are channel operations (\texttt{Probe}, \texttt{Receive}),
  which map naturally to LTS transitions.
\end{description}

%----------------------------------------------------------------------
\subsection{Challenges}
%----------------------------------------------------------------------

\begin{description}
\item[State explosion.]
  Even with the favorable blocking-point structure, realistic circuits
  have many processes.  The MiniMIPS reference design has 7~processes;
  a production pipeline might have 20--50.  With channel slack
  (buffering), each channel with slack~$d$ and width~$w$ adds up to
  $(2^w)^d$ states, though data abstraction can mitigate this.

\item[No specification language.]
  The Fulcrum CSP dialect has no mechanism for writing specifications or
  refinement assertions.  Deadlock freedom can be checked without a
  spec, but richer properties (e.g., ``the output channel produces
  values in sorted order'') require a specification language or
  assertion mechanism.

\item[Channel slack semantics.]
  FDR4 models channels as zero-buffer synchronous communications.
  The Intel toolchain supports non-unit channel slack (buffering),
  which must be modeled as explicit buffer processes in the LTS.
  This increases the state space but is straightforward to implement.

\item[Data-dependent control flow.]
  Some guards depend on data values (e.g., \texttt{x > 0}).
  Data-dependent branching requires either concrete enumeration of
  data values or symbolic techniques (SAT/SMT) to handle efficiently.

\item[Integration complexity.]
  FDR4 is a standalone tool written in Haskell with a C++ exploration
  engine.  Direct integration with the Modula-3/Scheme toolchain
  would require either a translation layer or a native implementation.
\end{description}

%======================================================================
\section{Proposed Approaches}
%======================================================================

We identify three approaches for adding formal verification to the
Fulcrum CSP toolchain, ranging from lowest implementation effort to
deepest integration.

%----------------------------------------------------------------------
\subsection{Approach~1: Translate to CSPM and Use FDR4}
%----------------------------------------------------------------------

Write a translator (in Scheme, within \texttt{cspc}) that converts
the compiled process network into CSPM syntax, emitting a
\texttt{.csp}\footnote{CSPM files conventionally use the \texttt{.csp}
extension; to avoid confusion with Intel \texttt{.csp} files, we would
use \texttt{.cspm} or \texttt{.fdr} as the extension.} file suitable
for FDR4.

\textbf{Advantages:}
\begin{itemize}
\item Leverages the mature, heavily optimized FDR4 engine.
\item Immediate access to all FDR4 analyses (deadlock, livelock,
  determinism, refinement).
\item Lowest implementation effort for basic functionality.
\end{itemize}

\textbf{Disadvantages:}
\begin{itemize}
\item FDR4 is proprietary software (University of Oxford); requires
  a license for commercial use.
\item Semantic gap: channel slack, data types, and probe operations
  require careful encoding.
\item Counterexample traces must be mapped back to the original
  Fulcrum CSP source, which adds complexity.
\item Maintenance burden: the translation must track changes in both
  the Fulcrum CSP dialect and CSPM.
\end{itemize}

%----------------------------------------------------------------------
\subsection{Approach~2: Native LTS Explorer}
%----------------------------------------------------------------------

Build a state-space exploration engine directly in the Fulcrum CSP
toolchain.  The engine would:

\begin{enumerate}
\item Compile each process to an LTS (states = blocking points,
  transitions = communication events).
\item Compose the per-process LTSs into a product LTS following the
  system topology declared in the \texttt{.sys} file.
\item Explore the product LTS using BFS/DFS with hash-based state
  storage, checking for deadlock (no enabled transitions) and
  optionally other properties.
\end{enumerate}

The natural implementation platform is Modula-3 (for performance)
with a Scheme API for scripting analyses, mirroring the existing
simulation architecture.

\textbf{Advantages:}
\begin{itemize}
\item No external tool dependencies; fully open-source.
\item Direct access to the AST and the simulation runtime's data
  structures (channels, process state).
\item Counterexample traces are in the native format.
\item Can exploit domain-specific optimizations: blocking-point
  reduction, data abstraction tailored to hardware channels.
\end{itemize}

\textbf{Disadvantages:}
\begin{itemize}
\item Significant implementation effort for the exploration engine.
\item Must implement compression techniques (partial-order reduction,
  symmetry reduction) to handle realistic circuits.
\item Initially limited to deadlock/livelock freedom; refinement
  checking requires substantially more infrastructure.
\end{itemize}

%----------------------------------------------------------------------
\subsection{Approach~3: SAT/BDD-Based Symbolic Model Checking}
%----------------------------------------------------------------------

Encode the process network as a Boolean transition relation and use
SAT solvers or BDD libraries for reachability analysis.  This is the
approach used by hardware model checkers (e.g., NuSMV, ABC).

\textbf{Advantages:}
\begin{itemize}
\item Handles large data widths efficiently (symbolic representation).
\item Mature tool ecosystem (MiniSat, CaDiCaL, CUDD).
\item Natural fit for hardware-oriented descriptions with fixed-width
  data.
\end{itemize}

\textbf{Disadvantages:}
\begin{itemize}
\item Largest implementation effort.
\item Encoding CSP concurrency semantics into Boolean relations is
  non-trivial and error-prone.
\item Loses the direct connection to CSP theory (traces, failures);
  properties must be expressed in CTL/LTL rather than as refinements.
\item Debugging counterexamples is harder (bit-level vs.\ process-level).
\end{itemize}

%======================================================================
\section{Recommended Approach}
%======================================================================

We recommend \textbf{Approach~2: a native LTS explorer}, starting
with deadlock freedom checking and incrementally adding capabilities.
The rationale is:

\begin{enumerate}
\item \textbf{No licensing constraints.}  FDR4's proprietary license
  is a significant obstacle for an open-source and Intel-internal
  toolchain.  A native implementation avoids this entirely.

\item \textbf{Domain-specific optimizations.}  The blocking-point
  structure of Fulcrum CSP processes means that each process's LTS is
  compact---typically tens to hundreds of states, not thousands.
  The main challenge is the product-space explosion, which can be
  addressed incrementally with partial-order reduction and
  compositional reasoning.

\item \textbf{Natural integration.}  The existing toolchain already
  compiles processes to state machines for simulation.  The same
  compilation pipeline can produce an LTS representation instead of
  (or in addition to) Modula-3 simulation code.  The Scheme scripting
  layer provides a natural interface for expressing verification
  queries.

\item \textbf{Incremental value.}  Deadlock freedom is the single
  most valuable property for asynchronous circuits, and it requires
  no specification language---only exhaustive exploration of the
  product state space.  This can be delivered as a first milestone
  with immediate practical value.

\item \textbf{Path to richer analyses.}  Once the LTS explorer
  exists, adding livelock detection (cycle detection on $\tau$-only
  paths) and trace refinement (product with a specification LTS) is
  a natural extension.
\end{enumerate}

The high-level architecture of the proposed verification engine is:

\begin{Verbatim}[samepage=true]
  .sys system description
           |
           v
    System graph (processes, channels, topology)
           |
           v
    Per-process LTS extraction
    (states = blocking points, transitions = chan ops)
           |
           v
    Product LTS composition (on-the-fly)
           |
           v
    +------+------+------+
    |      |      |      |
    v      v      v      v
  Deadlock  Livelock  Trace    Failures
  check     check     refine.  refine.
\end{Verbatim}

%======================================================================
\section{Project Plan}
%======================================================================

The project is organized into four phases, each delivering a usable
increment of verification capability.

%----------------------------------------------------------------------
\subsection{Phase~1: Per-Process LTS Extraction}
%----------------------------------------------------------------------

\textbf{Goal:} Extract an explicit LTS from a single compiled CSP
process.

\textbf{Deliverables:}
\begin{itemize}
\item A new compilation target in \texttt{cspc} that emits an LTS
  (set of states, initial state, set of labelled transitions) instead
  of Modula-3 simulation code.
\item States correspond to blocking points (channel send, receive,
  select).  Transitions are labelled with the channel operation
  (event).
\item Data variables are handled by bounded enumeration with
  configurable abstraction (e.g., collapse all values to a single
  abstract value for data-independent properties).
\item LTS output in a simple textual format (e.g., Aldebaran
  \texttt{.aut} or a custom S-expression format readable by
  \texttt{cspc}).
\item Scheme API: \texttt{(extract-lts \textit{process})} returns
  an LTS object.
\item Unit tests on the existing CSP test suite.
\end{itemize}

\textbf{Milestone:} Can extract and display the LTS of any single-process
CSP program (e.g., the producer and consumer from the \texttt{cspbuild}
examples).

%----------------------------------------------------------------------
\subsection{Phase~2: Product Composition and Deadlock Checking}
%----------------------------------------------------------------------

\textbf{Goal:} Compose per-process LTSs into a product system and
check for deadlock.

\textbf{Deliverables:}
\begin{itemize}
\item Product-state composition engine that synchronizes processes on
  shared channel events, following the topology from the \texttt{.sys}
  file.
\item Channel slack modeled as explicit FIFO buffer processes inserted
  between communicating processes.
\item On-the-fly BFS/DFS exploration with hash-based visited-state
  set.
\item Deadlock detection: a product state is a deadlock if no process
  can make progress (no enabled transitions in any component).
\item Counterexample generation: on deadlock, emit the trace
  (sequence of events) leading to the deadlocked state.
\item Scheme API: \texttt{(check-deadlock! \textit{sys-file})} returns
  either \texttt{\#t} (deadlock-free) or a counterexample trace.
\item Tested on the MiniMIPS reference design and the
  \texttt{cspbuild} examples.
\end{itemize}

\textbf{Milestone:} Can verify deadlock freedom of the producer-consumer
example and detect intentionally introduced deadlocks in test cases.

%----------------------------------------------------------------------
\subsection{Phase~3: State-Space Reduction}
%----------------------------------------------------------------------

\textbf{Goal:} Scale the explorer to handle realistic circuits
(20+ processes) through state-space reduction techniques.

\textbf{Deliverables:}
\begin{itemize}
\item \textbf{Partial-order reduction (POR):} Exploit the independence
  of events on disjoint channels to avoid exploring redundant
  interleavings.  The static topology from the \texttt{.sys} file
  directly provides the independence relation.
\item \textbf{Symmetry reduction:} Exploit structural symmetry in
  regular process arrays (e.g., pipeline stages, arbiters).
\item \textbf{Data abstraction:} Automatic abstraction of data
  channels for data-independent properties; user-guided abstraction
  for data-dependent properties.
\item \textbf{Compositional reasoning:} Verify subsystems independently
  and compose results, exploiting the hierarchical structure of
  asynchronous designs.
\item Performance benchmarks comparing exploration time and peak
  memory with and without each reduction.
\end{itemize}

\textbf{Milestone:} The MiniMIPS 7-process design is verified for
deadlock freedom within reasonable time and memory bounds (under
1~minute, under 1~GB).

%----------------------------------------------------------------------
\subsection{Phase~4: Livelock and Refinement Checking}
%----------------------------------------------------------------------

\textbf{Goal:} Extend the verification engine beyond deadlock freedom.

\textbf{Deliverables:}
\begin{itemize}
\item \textbf{Livelock detection:} Detect cycles in the product LTS
  that involve only internal ($\tau$) transitions and no visible
  progress.  Uses Tarjan's SCC algorithm on the $\tau$-reachable
  subgraph.
\item \textbf{Specification language:} A lightweight assertion
  mechanism, either embedded in \texttt{.sys} files or as a separate
  \texttt{.spec} file, for expressing trace properties (e.g.,
  ``channel~\texttt{out} always follows channel~\texttt{in}'').
\item \textbf{Trace refinement:} Check that the implementation's
  trace set is a subset of the specification's trace set, following
  the standard CSP refinement algorithm.
\item \textbf{Failures refinement} (stretch goal): Full
  failures-model refinement checking.
\item Documentation and user guide.
\end{itemize}

\textbf{Milestone:} Can express and verify a simple ordering property
on the producer-consumer example (``consumer receives values in the
order the producer sent them'').

%----------------------------------------------------------------------
\subsection{Timeline Summary}
%----------------------------------------------------------------------

\begin{center}
\begin{tabular}{llll}
\toprule
\textbf{Phase} & \textbf{Focus} & \textbf{Key Deliverable} & \textbf{Duration} \\
\midrule
1 & LTS extraction  & Per-process LTS from \texttt{cspc} & 6--8 weeks \\
2 & Deadlock checking & Product explorer + counterexamples & 8--10 weeks \\
3 & Reduction  & POR, symmetry, data abstraction & 8--12 weeks \\
4 & Livelock \& refinement & Spec language + refinement checker & 10--14 weeks \\
\bottomrule
\end{tabular}
\end{center}

Phases~1 and~2 together deliver the minimum viable product: a deadlock
checker for multi-process CSP systems.  Phases~3 and~4 extend
scalability and expressiveness incrementally.

%======================================================================
\section{Implementation Status}
%======================================================================

Phases~1 and~2 have been implemented in the \texttt{cspc} toolchain.
This section documents what was built, how it performs, and what
we learned.

%----------------------------------------------------------------------
\subsection{Phase~1: Per-Process LTS Extraction (Complete)}
%----------------------------------------------------------------------

LTS extraction is implemented as a Scheme pass within \texttt{cspc}.
The function \texttt{(extract-system-lts!~\textit{sys-file})} compiles
each process type referenced by a \texttt{.sys} file, runs the 9-pass
compiler to produce the intermediate representation, and then performs a
symbolic walk of the compiled state machine.  States correspond to
blocking points (channel sends and receives); transitions are labelled
with \texttt{(send~\textit{ch})} or \texttt{(recv~\textit{ch})} for
channel operations, and \texttt{tau} for internal (unobservable)
transitions such as guarded-repetition termination.

The per-process LTSs are compact.  For the dining philosophers
benchmark, the deterministic fork process has 5~states and 4~transitions;
the na\"ive philosopher has 5~states and 4~transitions.  This confirms
the blocking-point hypothesis: hardware CSP processes yield small LTSs.

%----------------------------------------------------------------------
\subsection{Phase~2a: Explicit-State Product Composition (Complete)}
%----------------------------------------------------------------------

The explicit-state deadlock checker composes per-process LTSs into a
product LTS on-the-fly using BFS.  Synchronization follows the
\texttt{.sys} topology: a \texttt{(send~\textit{ch})} transition in one
instance synchronizes with a \texttt{(recv~\textit{ch})} transition in
the connected instance; \texttt{tau} transitions are taken independently.
A product state is deadlocked if no component has an enabled transition.

When deadlock is found, the checker returns a counterexample trace
(sequence of synchronization events leading to the deadlocked state).
The BFS guarantees the trace is of minimum length.

The checker is implemented in Scheme (\texttt{product.scm}) and uses
the Modula-3 \texttt{RefSeq.T} type for the BFS queue.  It handles
all 15~acceptance tests, including $N = 2$ and $N = 3$ dining
philosophers with both deterministic and probe-based forks.

\textbf{Limitation:} The explicit-state approach enumerates product
states one at a time.  For $N = 8$ dining philosophers, the product
state space has $5^{16} \approx 1.5 \times 10^{11}$ states
(8~forks $\times$ 8~philosophers, each with 5~states), which is
infeasible for explicit enumeration.

%----------------------------------------------------------------------
\subsection{Phase~2b: BDD-Based Symbolic Deadlock Checking (Complete)}
%----------------------------------------------------------------------

To handle larger systems, a BDD-based symbolic deadlock checker was
implemented (\texttt{symbolic.scm}).  The checker encodes process states
as Boolean variables (2~bits per instance for a 5-state process),
builds the product transition relation $T(V, V')$ as a single BDD,
and computes reachability via a fixed-point iteration:
\[
  \mathit{Reached}_{i+1} = \mathit{Reached}_i \cup
    \mathit{Image}(\mathit{Reached}_i, T)
\]
where $\mathit{Image}$ conjoins the reached set with $T$, existentially
quantifies out the current-state variables, and renames next-state
variables to current-state variables.

After convergence, deadlocked states are those in the reached set with
no successor: $\mathit{Deadlock} = \mathit{Reached} \wedge \neg \mathit{HasSucc}$,
where $\mathit{HasSucc} = \exists V'. T(V, V')$.

The implementation uses the existing \texttt{bdd} library (part of
the CM3 tree).  All BDD operations are accessible from Scheme via
auto-generated stubs (\texttt{SchemeStubs}).

\textbf{Results.}
The symbolic checker handles $N = 8$ dining philosophers (16~instances,
32~BDD variables) and reaches the fixed point in 8~iterations.
The following table shows the BDD size of the reached set at each
iteration:

\begin{center}
\begin{tabular}{rr}
\toprule
\textbf{Iteration} & \textbf{BDD nodes} \\
\midrule
0 & 194 \\
1 & 474 \\
2 & 866 \\
3 & 1,230 \\
4 & 1,454 \\
5 & 1,542 \\
6 & 1,562 \\
7 & 1,564 \\
8 (fixed point) & 1,564 \\
\bottomrule
\end{tabular}
\end{center}

The BDD representation is compact: 1,564~nodes encode the full
reachable state set (which contains an exponential number of concrete
states).  The deadlock BDD has 66~nodes.  Execution time on Apple~M2
is approximately 35~seconds, dominated by the image computation at each
iteration.  Peak memory is approximately 2.7~GB.

%----------------------------------------------------------------------
\subsection{Memory Analysis and BDD Garbage Collection}
%----------------------------------------------------------------------

The 2.7~GB peak memory for $N = 8$ dining philosophers motivated an
investigation into BDD garbage collection.  The \texttt{bdd} library
maintains two data structures per variable that grow without bound:

\begin{description}
\item[Unique tables.]  Map $(l, r) \to \mathit{node}$, ensuring
  that structurally identical BDD nodes are pointer-equal (canonicity).
  Every BDD node ever created under a variable remains in this table.
\item[Operation caches.]  Map $(\mathit{input}_1, \mathit{input}_2)
  \to \mathit{result}$ for each operation (And, Or, Not, MakeTrue,
  MakeFalse).  These provide memoization for recursive BDD operations.
\end{description}

Both tables hold strong references to BDD nodes, preventing the
Modula-3 garbage collector from reclaiming them even when no user
code references the intermediate results.

\textbf{Solution: explicit mark-and-sweep.}
Four procedures were added to the \texttt{BDD} interface
(in the CM3 repository):

\begin{description}
\item[\texttt{ClearCaches()}.]  Drops all operation caches.
  $O(\text{number of variables})$, no DAG traversal.
\item[\texttt{MarkRoot(t)}.]  Adds \texttt{t} to a GC root set.
\item[\texttt{UnmarkRoots()}.]  Clears the GC root set.
\item[\texttt{CollectGarbage()}.]  Marks all nodes reachable from the
  root set, rebuilds each variable's unique table keeping only live
  entries, clears all operation caches, and clears the root set.
\end{description}

These are callable from Scheme via auto-generated stubs
(\texttt{BDD.ClearCaches}, \texttt{BDD.MarkRoot}, etc.).

\textbf{Benchmark results.}
Inserting \texttt{CollectGarbage} between fixed-point iterations on
the $N = 8$ dining philosophers benchmark yielded the following:

\begin{center}
\begin{tabular}{lrr}
\toprule
& \textbf{Without GC} & \textbf{With GC} \\
\midrule
Peak RSS     & 2,702~MB & 2,749~MB \\
Wall time    & 34.7~s   & 47.8~s \\
Result       & DEADLOCK & DEADLOCK \\
\bottomrule
\end{tabular}
\end{center}

The GC does not reduce peak memory and adds approximately 38\%
overhead.  The reason is that peak memory is dominated by
\emph{within-iteration} temporaries: the image computation
$\texttt{BDD.And}(\mathit{reached}, T)$ creates a large intermediate
BDD that is live throughout the existential quantification loop.
The GC runs \emph{between} iterations, by which time the Modula-3
runtime has already grown its heap to accommodate the peak, and
does not return memory to the operating system.

\textbf{Conclusion.}
The GC hooks are retained in the BDD library for future use but
are not called by the symbolic checker.  Reducing memory for
large problems will require algorithmic improvements:
partitioned transition relations (conjunctive/disjunctive
partitioning of $T$), saturation-based image computation, or
BDD variable reordering.  These are natural Phase~3 targets.

%======================================================================
\section{References}
%======================================================================

\begin{thebibliography}{99}

\bibitem{hoare1985}
C.~A.~R. Hoare,
\emph{Communicating Sequential Processes},
Prentice Hall, 1985.
Available online: \url{http://www.usingcsp.com/}

\bibitem{roscoe2010}
A.~W. Roscoe,
\emph{Understanding Concurrent Systems},
Springer, 2010.

\bibitem{fdr4}
T.~Gibson-Robinson, P.~Armstrong, A.~Boulgakov, A.~W. Roscoe,
``FDR3: A Modern Refinement Checker for CSP,''
\emph{Proc.\ TACAS}, LNCS 8413, pp.~187--201, 2014.
FDR4 documentation: \url{https://cocotec.io/fdr/}

\bibitem{martin1990}
A.~J. Martin,
``Programming in VLSI: From Communicating Processes to Delay-Insensitive Circuits,''
\emph{Developments in Concurrency and Communication}, UT Year of Programming Series,
Addison-Wesley, pp.~1--64, 1990.

\bibitem{martin2006}
A.~J. Martin, M.~Nystr\"om,
``Asynchronous Techniques for System-on-Chip Design,''
\emph{Proc.\ IEEE}, vol.~94, no.~6, pp.~1089--1120, June 2006.

\bibitem{josephs1995}
M.~B. Josephs, J.~T. Udding,
``An Overview of DI Algebra,''
\emph{Proc.\ Hawaii International Conference on System Sciences}, 1995.

\bibitem{por}
P.~Godefroid,
\emph{Partial-Order Methods for the Verification of Concurrent Systems},
LNCS 1032, Springer, 1996.

\bibitem{mcmillan1993}
K.~L. McMillan,
\emph{Symbolic Model Checking},
Kluwer Academic Publishers, 1993.

\bibitem{supercomb}
A.~W. Roscoe,
``Compiling Shared-Nothing CSP to Supercombinators: A Progress Report,''
Technical report, University of Oxford, 2013.

\end{thebibliography}

\end{document}
